% 核心结论：同步转速与滑差
% 图原型：纯几何／对比。物理网络按本题节点、端口和符号构建。
\documentclass[tikz,border=7pt]{standalone}
\usepackage{fontspec,xeCJK,amsmath,amssymb}
\IfFontExistsTF{Times New Roman}{\setmainfont{Times New Roman}}{\setmainfont{TeX Gyre Termes}}
\IfFontExistsTF{Heiti SC}{\setCJKmainfont[AutoFakeBold=1.5]{Heiti SC}}{\setCJKmainfont[AutoFakeBold=1.5]{FandolHei-Regular}}
\usetikzlibrary{arrows.meta,calc,angles,quotes}
\definecolor{ink}{HTML}{243745}
\definecolor{blue}{HTML}{275DAB}
\definecolor{green}{HTML}{177C63}
\definecolor{red}{HTML}{BD4935}
\definecolor{purple}{HTML}{715393}
\begin{document}
\begin{tikzpicture}[font=\fontsize{9.5}{12}\selectfont,text=ink,
 v/.style={line width=1.1pt,-{Stealth[length=2.4mm,width=1.5mm]}},
 axis/.style={ink!65,line width=.45pt,-{Stealth[length=1.7mm,width=1.1mm]}},
 aux/.style={ink!40,densely dashed,line width=.5pt},
 lab/.style={fill=white,inner sep=2pt},
 card/.style={draw=blue!25,fill=blue!3,rounded corners=3pt,text width=6cm,align=left,inner sep=10pt}]
\path[use as bounding box] (-8,-3.35) rectangle (8,3.65);
\node[font=\fontsize{10.5}{13}\selectfont\bfseries,text width=15cm,align=center] at (0.00000,3.22000) {M6-077　同步转速与滑差};
\begin{scope}[shift={(-4,0)}]
\draw[ink!60,fill=blue!3,thick] (0,0) circle(1.7);
\draw[ink,fill=white,thick] (0,0) circle(.85);
\node[] at (0.00000,0.00000) {转子};
\draw[v,blue] (1.7,0) arc[start angle=0,end angle=140,radius=1.7];
\node[text=blue] at (0.15000,2.15000) {$n_s$};
\draw[v,green] (.85,0) arc[start angle=0,end angle=130,radius=.85];
\node[text=green] at (1.35000,0.70000) {$n_r$};
\node[] at (0.00000,-2.40000) {正常电动状态：$0<n_r<n_s$。};
\end{scope}
\begin{scope}[shift={(4,0)}]
\node[card] at (0.00000,1.15000) {$n_s=1500\;r/min,\quad n_r=1440$};
\end{scope}
\begin{scope}[shift={(4,0)}]
\node[card] at (0.00000,-1.15000) {$s=.0400,\quad f_r=2Hz$};
\end{scope}
\node[text width=15cm,align=center] at (0.00000,-3.00000) {同步磁场比转子快，转子电流频率由相对转速确定；交换两相先改变磁场方向，转子需经历动态过程。};

\end{tikzpicture}\end{document}